\section{Related Work}
\label{sec:related}
\paragraph{Multi-task benchmarks and their aggregation.}
Multi-subject and multi-metric suites such as MMLU~\citep{hendrycks2020mmlu}, BIG-Bench Hard~\citep{suzgun2022bbh} and HELM~\citep{liang2022helm} report one score per sub-dataset, and public leaderboards then collapse these scores into a single ranking, most often with an (optionally normalised) arithmetic mean.
\citet{dehghani2021lottery} show that task selection and other benchmark design choices can change which method looks best; our question is complementary: for a \emph{fixed} task set, how robust is the first place to a developer's \emph{action}, and how does that depend on the aggregation rule?

\paragraph{Aggregation as social choice.}
Ranking models from several score columns is a voting problem, with tasks as voters and models as candidates~\citep{brandt2016handbook}. \citet{lanctot2023socialchoice} advocate social-choice rules for evaluating agents over many tasks. \citet{rofin2023votenrank} apply such rules to re-rank NLP benchmark results and compare the rankings; they do not ask what a model must change to be first. We use the lens to pick rules (Borda, Condorcet-winner requirement, plurality) alongside score-based rules (mean, power means, trimmed mean, median, min), and we ask a question social choice calls \emph{constructive manipulation} or bribery: the cheapest change of one candidate's profile that makes it win~\citep{bartholdi1989manipulating,faliszewski2016control}, or bribery \citep{faliszewski2009bribery}. That literature mostly studies the computational complexity of changing ballots; here the changed object is the target's own real-valued scores, within bounds, the score rules reduce to covering problems with closed forms (Propositions~\ref{prop:mean}--\ref{prop:median}), and the content is how the cost behaves on real score matrices. The Open LLM Leaderboard documents a per-task normalisation of raw scores before averaging~\citep{fourrier2024ollb}; we take its normalised scores and mean as the leaderboard's own rule.

\paragraph{Gaming, leakage and contamination.}
\citet{singh2025illusion} document distortions of the Chatbot Arena leaderboard (private testing of many variants with selective disclosure, unequal data access). \citet{zhou2023cheater} show that benchmark leakage into training can dramatically boost evaluation results, and \citet{yang2023rephrased} show that rephrased or translated test samples evade n-gram decontamination yet let a model overfit a benchmark. These works establish \emph{that} targeted training on benchmark-related data inflates scores. We take that effect as an assumption (a developer may raise scores on chosen sub-datasets, up to a stated cap) and ask how much of it a given aggregation rule requires to reach first place.
\citet{chen2026incentives} model benchmarking as a Stackelberg game between a designer and developers who allocate post-training resources to benchmark-specific gains, and prove equilibrium results. Their central quantity, the ``just-overtake'' effort, is the least extra training a model needs to pass the model ranked directly above it on one benchmark; they estimate it from fitted training trajectories in a case study. Their model predicts how much post-training effort developers spend in equilibrium on one benchmark; it has no notion of which sub-dataset to touch or of the aggregation rule. We study \emph{multi-task} leaderboards, ask what it costs to reach \emph{first} place, and compare \emph{aggregation rules} on real public score matrices, with the cost measured in sub-datasets touched and score points rather than training steps; other models are held fixed and no equilibrium is computed.
